\ifthenelse{\equal{\lang}{ko}}{%
 법률, 규제, 정부 정보는 시민의 권리와 의무, 공공 서비스 이용, 사회·경제적 참여의 근간을 이룬다. 그러나 이러한 정보는 복잡한 전문 용어, 파편화된 출처, 그리고 언어 장벽으로 인해 비전문가와 다국어 공동체에게 여전히 접근하기 어려운 자원으로 남아 있다 \cite{zhong2020legalai, chalkidis2020legalbert}. 이민자, 소상공인, 취약 계층, 그리고 자국어가 공식 문서의 언어와 다른 사람들에게 이러한 장벽은 정보 격차를 심화시키고 공정한 공공 서비스 접근을 저해한다 \cite{aletras2016echr}.

여러 관할권에서 법령, 판례, 계약, 행정 공지는 서로 다른 언어와 형식으로 지속적으로 갱신된다. 유럽연합의 AI 법안 \cite{EUAIAct2024}과 같은 대규모 규제 프레임워크는 다국어로 공표되지만, 일반 시민이 이를 자신의 언어로 정확히 이해하고 자신의 상황에 적용하기는 쉽지 않다. 기존의 접근 방식—키워드 검색, 수작업 번역, 정적 정보 포털—은 의미적 깊이가 부족하거나, 최신성을 유지하지 못하거나, 저자원 언어를 충분히 다루지 못한다 \cite{niklaus2023multilegal}.

본 논문에서는 다양한 관할권과 언어에 걸쳐 거버넌스 관련 정보를 \textbf{수집·분류·분석·요약}하도록 설계된 다국어 AI 프레임워크인 \textbf{LexSense}를 제안한다. LexSense는 (1) 자동 변화 탐지, (2) 다국어 자연어 처리 \cite{conneau2020xlmr, xue2021mt5}, (3) 개념 드리프트(concept drift) 모니터링 \cite{gama2014drift}, (4) 근거 기반(evidence-grounded) 분석 \cite{lewis2020rag}, (5) 대규모 언어 모델(LLM) 기반 보고 \cite{openai2023gpt4, touvron2023llama2}를 하나의 엔드투엔드 파이프라인으로 통합한다. 핵심 목표는 규제 준수를 대신하는 것이 아니라, 법률 및 정부 지식에 대한 \emph{포용적 접근}을 지원하는 것이다.
~\\
본 연구의 기여는 세 가지다.
\textbf{첫째, 다국어 접근성 프레임워크:} 변화 탐지, 다국어 분류, 근거 기반 요약, LLM 보고를 결합한 모듈형 프레임워크를 설계·구현하여, 비전문가와 다국어 사용자가 공공 정보에 접근하는 데 따르는 언어적·기술적 장벽을 낮춘다 \cite{devlin2019bert, reimers2020multilingual}.
\textbf{둘째, 다국어 벤치마크(GovSense-1k):} 근거 주석이 포함된 1,200건의 거버넌스·계약·소송·자산 관련 레코드로 구성된 다국어 벤치마크를 공개하고, 이를 영어·한국어·프랑스어에 걸쳐 검증하였다 \cite{chalkidis2022lexglue, hendrycks2021cuad}.
\textbf{셋째, 포괄적 평가:} 영어·한국어·프랑스어에 대한 주 평가와 스페인어·아랍어·베트남어에 대한 스트레스 테스트를 수행하여, 다국어·근거 기반 AI가 공공 정보 접근성을 향상시킬 수 있음을 정량적으로 입증하였다 \cite{nllb2022, hu2020xtreme}.

이러한 기술적 기여를 넘어, 본 연구는 책임 있는 배포를 위한 사실성(factuality), 투명성, 인간 감독(human oversight) 요구사항을 함께 논의한다 \cite{ji2023hallucination, maynez2020faithfulness, weidinger2021taxonomy}. 종합하면, LexSense는 자연어 처리, 법·정책 정보학, 책임 있는 AI의 교차점에서, 공공 지식에 대한 \textbf{포용적 접근}을 위한 실용적이면서도 확장 가능한 인프라를 제시한다 \cite{bommasani2021foundation}.



}{%
% -------------------- English Version --------------------
Legal, regulatory, and government information underpins citizens' rights and obligations, access to public services, and socioeconomic participation. Yet such information remains difficult to access for non-experts and multilingual communities because of complex terminology, fragmented sources, and language barriers \cite{zhong2020legalai, chalkidis2020legalbert}. For immigrants, small business owners, vulnerable populations, and people whose native language differs from that of official documents, these barriers widen information gaps and hinder equitable access to public services \cite{aletras2016echr}.

Across jurisdictions, statutes, case law, contracts, and administrative notices are continuously updated in different languages and formats. Large regulatory frameworks such as the European Union AI Act \cite{EUAIAct2024} are published in multiple languages, yet ordinary citizens still struggle to understand them accurately in their own language and to apply them to their own circumstances. Existing approaches---keyword search, manual translation, and static information portals---lack semantic depth, fail to stay current, or provide insufficient coverage of low-resource languages \cite{niklaus2023multilegal}.

In this paper, we introduce \textbf{LexSense}, a multilingual AI framework designed to \textbf{collect, classify, analyze, and summarize} governance-related information across different jurisdictions and languages. LexSense integrates (1) automated change detection, (2) multilingual natural language processing \cite{conneau2020xlmr, xue2021mt5}, (3) concept-drift monitoring \cite{gama2014drift}, (4) evidence-grounded analysis \cite{lewis2020rag}, and (5) large language model (LLM)-based reporting \cite{openai2023gpt4, touvron2023llama2} into a single end-to-end pipeline. Its central goal is not to replace regulatory compliance functions but to support \emph{inclusive access} to legal and governmental knowledge.
~\\
Our contributions are threefold.
\textbf{First, a multilingual accessibility framework:} we design and implement a modular framework that combines change detection, multilingual classification, evidence-grounded summarization, and LLM reporting, lowering the linguistic and technical barriers that non-experts and multilingual users face when accessing public information \cite{devlin2019bert, reimers2020multilingual}.
\textbf{Second, a multilingual benchmark (GovSense-1k):} we release a multilingual benchmark of 1,200 governance, contract, lawsuit, and asset-related records with evidence annotations, validated across English, Korean, and French \cite{chalkidis2022lexglue, hendrycks2021cuad}.
\textbf{Third, comprehensive evaluation:} we conduct a main evaluation on English, Korean, and French, together with stress tests on Spanish, Arabic, and Vietnamese, quantitatively demonstrating that multilingual, evidence-grounded AI can improve access to public information \cite{nllb2022, hu2020xtreme}.

Beyond these technical contributions, we discuss factuality, transparency, and human oversight requirements for responsible deployment \cite{ji2023hallucination, maynez2020faithfulness, weidinger2021taxonomy}. Taken together, LexSense offers a practical yet scalable infrastructure for \textbf{inclusive access} to public knowledge, at the intersection of natural language processing, legal and policy informatics, and responsible AI \cite{bommasani2021foundation}.

}
